% Ementa — Núcleo de Análise Quantitativa
% Liga de Mercado Financeiro — UFU | 2026.1
% Diretoria de Pesquisa — Arthur Pires Lopes

\documentclass[12pt,a4paper]{article}

% Pacotes necessários
\usepackage[utf8]{inputenc}
\usepackage[T1]{fontenc}
\usepackage[portuguese]{babel}
\usepackage{graphicx}
\usepackage{geometry}
\usepackage{titlesec}
\usepackage{enumitem}
\usepackage{xcolor}
\usepackage{hyperref}
\usepackage{fancyhdr}
\usepackage{setspace}
\usepackage{parskip}
\usepackage{eso-pic}
\usepackage{tikz}
\usepackage{wrapfig}
\usepackage{booktabs}
\usepackage{tabularx}
\usepackage{longtable}
\usepackage{array}
\usepackage{colortbl}
\usepackage{fontawesome5}
\usetikzlibrary{positioning}

% Configuração de margens
\geometry{left=2.5cm, right=2.5cm, top=3cm, bottom=2.5cm}

% Cores personalizadas (padrão LMF UFU)
\definecolor{ufuazul}{RGB}{0, 51, 102}
\definecolor{cinzaescuro}{RGB}{64, 64, 64}
\definecolor{azulclaro}{RGB}{220, 232, 245}
\definecolor{verdeok}{RGB}{34, 139, 34}

% Configuração de hyperlinks
\hypersetup{
    colorlinks=true,
    linkcolor=ufuazul,
    urlcolor=ufuazul
}

% Configuração de títulos
\titleformat{\section}{\Large\bfseries\color{ufuazul}}{\thesection.}{0.5em}{}
\titleformat{\subsection}{\large\bfseries\color{cinzaescuro}}{\thesubsection.}{0.5em}{}

% Espaçamento
\setstretch{1.15}

% ============================================
% COMANDO PARA IMAGEM DE FUNDO NA CAPA
% ============================================
\newcommand{\BackgroundPic}{%
    \begin{tikzpicture}[remember picture, overlay]
        \node[opacity=0.15, inner sep=0pt, anchor=center] at (current page.center) {
            \includegraphics[width=\paperwidth, height=\paperheight]{fundo.png}
        };
    \end{tikzpicture}
}

% Cabeçalho e rodapé
\pagestyle{fancy}
\fancyhf{}
\fancyhead[L]{%
    \raisebox{-0.3\height}{\includegraphics[height=1.2cm]{logo_lmf.png}}%
}
\fancyhead[R]{\small\textcolor{gray}{Núcleo Quant | Ementa 2026.1}}
\fancyfoot[C]{\thepage}
\renewcommand{\headrulewidth}{0pt}

% Estilo para capa
\fancypagestyle{capa}{
    \fancyhf{}
    \renewcommand{\headrulewidth}{0pt}
    \renewcommand{\footrulewidth}{0pt}
}

% Novo tipo de coluna para tabela
\newcolumntype{C}[1]{>{\centering\arraybackslash}m{#1}}
\newcolumntype{L}[1]{>{\raggedright\arraybackslash}m{#1}}

\begin{document}

% ============================================
% CAPA
% ============================================
\thispagestyle{capa}

\begin{tikzpicture}[remember picture, overlay]
    \node[opacity=0.55, inner sep=0pt] at (current page.center) {
        \includegraphics[width=\paperwidth, height=\paperheight]{fundo.png}
    };
\end{tikzpicture}

\begin{center}
    \vspace*{1cm}
    
    \includegraphics[height=3cm]{logo_lmf.png}
    
    \vspace{1.5cm}
    
    {\Large\textbf{LIGA DE MERCADO FINANCEIRO --- UFU}}\\[0.3cm]
    {\large Diretoria de Pesquisa — Ementa 2026.1}\\[1.5cm]
    
    \rule{0.8\textwidth}{1pt}\\[0.5cm]
    {\LARGE\textbf{Núcleo de Análise Quantitativa}}\\[0.3cm]
    {\large\textit{Quant}}\\[0.5cm]
    \rule{0.8\textwidth}{1pt}
    
    \vspace{2cm}
    
    {\Large\textbf{Arthur Pires Lopes}}\\[0.3cm]
    {\large Diretor de Pesquisa}\\[0.3cm]
    {\large Engenharia Elétrica — 6° Período}
    
    \vspace{1.5cm}
    
    {\normalsize\textcolor{cinzaescuro}{Challenge de referência: \textbf{Itaú Quant Challenge}}}
    
    \vfill
    
    \includegraphics[height=1.5cm]{logo_lmf.png}
    
    \vspace{1cm}
\end{center}

\newpage

% ============================================
% CONTEÚDO
% ============================================

% ============================================
\section{Objetivo do Núcleo}
% ============================================

O Núcleo de Análise Quantitativa tem como missão \textbf{formar membros em métodos quantitativos aplicados a finanças}, desenvolvendo competências em programação, estatística, modelagem e gestão de risco que se traduzam em \textbf{performance real em challenges} — em especial o Itaú Quant Challenge.

O núcleo preenche uma lacuna na Liga: atualmente não há espaço estruturado para quem quer trabalhar com \textbf{dados, modelos e automação} no contexto de mercados. A proposta é criar esse espaço com processo claro, entregas aplicadas e padrão de qualidade reproduzível.

\vspace{0.3cm}
\noindent\textbf{Ao final do ciclo, o membro será capaz de:}
\begin{itemize}[leftmargin=*]
    \item Coletar, tratar e analisar dados financeiros com Python;
    \item Construir e validar estratégias quantitativas com backtesting rigoroso;
    \item Mensurar e gerenciar risco de portfólio (VaR, CVaR, drawdown, stress test);
    \item Aplicar técnicas de Machine Learning a problemas financeiros;
    \item Desenvolver e operar bots de trading algorítmico;
    \item Comunicar resultados com clareza técnica (relatórios e notebooks documentados).
\end{itemize}

% ============================================
\section{Pré-requisitos e Ferramentas}
% ============================================

\subsection{Pré-requisitos recomendados}
\begin{itemize}[leftmargin=*]
    \item Noções básicas de programação (qualquer linguagem; Python é preferencial);
    \item Estatística descritiva (média, variância, distribuições);
    \item Interesse genuíno por mercado financeiro e métodos quantitativos.
\end{itemize}

\textit{Observação:} não é necessário dominar todos os pré-requisitos — as primeiras aulas nivelam a base. O mais importante é disposição para estudar e praticar.

\subsection{Stack Técnico}
\begin{center}
\renewcommand{\arraystretch}{1.4}
\begin{tabular}{l l}
    \toprule
    \textbf{Ferramenta} & \textbf{Uso} \\
    \midrule
    \textbf{Python 3.10+} & Linguagem principal \\
    \textbf{Jupyter Notebook / VS Code} & Ambiente de desenvolvimento \\
    \textbf{pandas / NumPy} & Manipulação e cálculo numérico \\
    \textbf{Matplotlib / Plotly} & Visualização de dados \\
    \textbf{SciPy / statsmodels} & Estatística e séries temporais \\
    \textbf{scikit-learn} & Machine Learning \\
    \textbf{backtrader / bt / vectorbt} & Backtesting de estratégias \\
    \textbf{yfinance / APIs} & Coleta de dados de mercado \\
    \textbf{Git + GitHub} & Versionamento e repositório \\
    \bottomrule
\end{tabular}
\end{center}

% ============================================
\section{Cronograma de Aulas — 12 Semanas}
% ============================================

O ciclo é composto por \textbf{12 encontros quinzenais}, alternando teoria e prática. Cada aula tem duração estimada de 1h30 a 2h. O cronograma está organizado em três blocos progressivos:

\vspace{0.3cm}

\noindent
\textbf{\textcolor{ufuazul}{\faIcon{layer-group} Bloco 1 — Fundamentos (Semanas 1--4)}}

\vspace{0.3cm}
\renewcommand{\arraystretch}{1.5}
\begin{tabularx}{\textwidth}{C{1.8cm} L{3.5cm} X}
    \toprule
    \rowcolor{azulclaro}
    \textbf{Semana} & \textbf{Tema} & \textbf{Conteúdo} \\
    \midrule
    
    \textbf{1} & Python para Finanças I &
    Ambiente de trabalho (Jupyter, VS Code, Git). Introdução a \texttt{pandas} e \texttt{NumPy}. Importação e manipulação de séries de preços. Cálculo de retornos simples e log-retornos. \\
    
    \textbf{2} & Python para Finanças II &
    Coleta de dados via \texttt{yfinance} e APIs. Tratamento de dados faltantes. Visualização de séries temporais com \texttt{Matplotlib} e \texttt{Plotly}. Análise exploratória de ativos brasileiros e internacionais. \\
    
    \textbf{3} & Estatística Aplicada &
    Estatística descritiva de retornos. Distribuições (normal, t-Student, caudas pesadas). Testes de hipótese (normalidade, estacionariedade). Correlação e covariância entre ativos. \\
    
    \textbf{4} & Séries Temporais &
    Conceito de estacionariedade e testes ADF. Autocorrelação (ACF/PACF). Modelos ARIMA. Modelagem de volatilidade: GARCH. Aplicação prática em dados reais. \\
    
    \bottomrule
\end{tabularx}

\vspace{0.8cm}

\noindent
\textbf{\textcolor{ufuazul}{\faIcon{shield-alt} Bloco 2 — Risco e Backtesting (Semanas 5--8)}}

\vspace{0.3cm}
\renewcommand{\arraystretch}{1.5}
\begin{tabularx}{\textwidth}{C{1.8cm} L{3.5cm} X}
    \toprule
    \rowcolor{azulclaro}
    \textbf{Semana} & \textbf{Tema} & \textbf{Conteúdo} \\
    \midrule
    
    \textbf{5} & Gestão de Risco I &
    Value at Risk (VaR): paramétrico, histórico e por simulação de Monte Carlo. Conditional VaR (CVaR / Expected Shortfall). Drawdown máximo e métricas de cauda. \\
    
    \textbf{6} & Gestão de Risco II &
    Teoria de portfólio de Markowitz. Fronteira eficiente e otimização média-variância. Diversificação e matriz de correlação. Stress testing e análise de cenários. \\
    
    \textbf{7} & Backtesting I &
    Fundamentos de backtesting. Métricas de performance: Sharpe, Sortino, Calmar, Treynor. Vieses comuns: \textit{look-ahead bias}, \textit{survivorship bias}, \textit{overfitting}. Introdução ao framework \texttt{backtrader}. \\
    
    \textbf{8} & Backtesting II &
    Implementação prática de estratégias (momentum, mean reversion, breakout). Walk-forward analysis e validação out-of-sample. Relatório de backtest padronizado. \textbf{Entrega intermediária.} \\
    
    \bottomrule
\end{tabularx}

\vspace{0.8cm}

\noindent
\textbf{\textcolor{ufuazul}{\faIcon{robot} Bloco 3 — Machine Learning, Bots e Projeto Final (Semanas 9--12)}}

\vspace{0.3cm}
\renewcommand{\arraystretch}{1.5}
\begin{tabularx}{\textwidth}{C{1.8cm} L{3.5cm} X}
    \toprule
    \rowcolor{azulclaro}
    \textbf{Semana} & \textbf{Tema} & \textbf{Conteúdo} \\
    \midrule
    
    \textbf{9} & Machine Learning I &
    Regressão linear e logística aplicadas a finanças. Feature engineering com dados de mercado (retornos defasados, volatilidade, indicadores técnicos). Validação temporal (\textit{time-series split}). \\
    
    \textbf{10} & Machine Learning II &
    Modelos ensemble: Random Forest, Gradient Boosting (XGBoost / LightGBM). Seleção de features e regularização. Overfitting em finanças: por que modelos complexos falham. Interpretação de resultados. \\
    
    \textbf{11} & Algoritmos de Trading &
    Lógica de bots de trading. Gestão de ordens (limit, market, stop). Simulação e paper trading. Integração com APIs de corretoras. Custos de transação e slippage. \\
    
    \textbf{12} & Projeto Final &
    Apresentação da estratégia completa: tese $\rightarrow$ modelo $\rightarrow$ backtest $\rightarrow$ resultado. Cada membro/grupo defende sua abordagem. Feedback coletivo e documentação para o repositório da Liga. \\
    
    \bottomrule
\end{tabularx}

% ============================================
\section{Metodologia}
% ============================================

\begin{itemize}[leftmargin=*]
    \item \textbf{Formato:} aulas quinzenais teórico-práticas (alternando com os demais núcleos da Diretoria de Pesquisa), com duração de 1h30 a 2h.
    
    \item \textbf{Dinâmica:} cada encontro combina exposição conceitual curta ($\sim$30 min) com exercício \textit{hands-on} em Jupyter Notebook ($\sim$60 min). O foco é \textbf{aprender fazendo}.
    
    \item \textbf{Mini-entregas:} a cada 2--4 semanas, os membros submetem um notebook documentado com a aplicação prática do bloco. Isso garante cadência e evita acúmulo no final.
    
    \item \textbf{Discussão de referências:} ao longo do ciclo, papers e capítulos de livros são indicados como leitura complementar e discutidos brevemente no início das aulas.
    
    \item \textbf{Repositório compartilhado:} todo material (slides, notebooks, datasets, templates) ficará versionado no GitHub da Liga, criando acervo reaproveitável para os próximos semestres.
    
    \item \textbf{Suporte assíncrono:} grupo dedicado para dúvidas, compartilhamento de materiais e discussões entre encontros.
\end{itemize}

% ============================================
\section{Entregas e Avaliação}
% ============================================

\subsection{Entregas esperadas}

\renewcommand{\arraystretch}{1.4}
\begin{tabularx}{\textwidth}{l L{3.5cm} X c}
    \toprule
    \textbf{\#} & \textbf{Entrega} & \textbf{Descrição} & \textbf{Semana} \\
    \midrule
    1 & Análise Exploratória & Notebook com coleta, tratamento e visualização de dados de pelo menos 5 ativos. Estatísticas descritivas e correlação. & 4 \\
    \addlinespace
    2 & Relatório de Risco & Cálculo de VaR, CVaR e drawdown para um portfólio. Otimização de Markowitz com fronteira eficiente. & 6 \\
    \addlinespace
    3 & Backtest Completo & Implementação de uma estratégia quantitativa com backtest rigoroso, métricas de performance e análise de robustez. & 8 \\
    \addlinespace
    4 & Projeto Final & Estratégia completa (tese + modelo + backtest + resultado), apresentada e documentada. & 12 \\
    \bottomrule
\end{tabularx}

\subsection{Critérios de avaliação (rubrica)}

Cada entrega será avaliada em 4 dimensões com peso igual:

\begin{enumerate}[leftmargin=*]
    \item \textbf{Rigor técnico:} os métodos estão corretos? Os pressupostos foram verificados? O código é limpo e funcional?
    
    \item \textbf{Clareza:} o notebook é auto-explicativo? Há narrativa conectando as etapas? Gráficos estão rotulados?
    
    \item \textbf{Reprodutibilidade:} outra pessoa consegue rodar o notebook do zero e obter os mesmos resultados?
    
    \item \textbf{Interpretação:} os resultados foram discutidos de forma crítica? Limitações foram apontadas? Há conclusão objetiva?
\end{enumerate}

\textbf{Nota:} a meta não é perfeição — é \textit{progresso consistente}. Entregar no prazo com estrutura razoável vale mais do que entregar atrasado com código impecável.

% ============================================
\section{Referências e Recursos}
% ============================================

\subsection{Livros recomendados}

\begin{enumerate}[leftmargin=*]
    \item \textbf{Marcos López de Prado} — \textit{Advances in Financial Machine Learning} (Wiley, 2018).\\
    Referência central em ML aplicado a finanças. Cobre backtesting, feature engineering, overfitting e meta-labeling.
    
    \item \textbf{Ernest P. Chan} — \textit{Quantitative Trading: How to Build Your Own Algorithmic Trading Business} (Wiley, 2nd ed.).\\
    Guia prático para construção de estratégias quantitativas, gestão de risco e operação de bots.
    
    \item \textbf{Yves Hilpisch} — \textit{Python for Finance: Mastering Data-Driven Finance} (O'Reilly, 2018).\\
    Cobertura completa de Python aplicado a dados financeiros, derivativos e simulação.
    
    \item \textbf{Stefan Jansen} — \textit{Machine Learning for Algorithmic Trading} (Packt, 2nd ed., 2020).\\
    Foco em pipelines completos de ML para trading, com código reproduzível em Python.
    
    \item \textbf{Andrew Ang} — \textit{Asset Management: A Systematic Approach to Factor Investing} (Oxford, 2014).\\
    Base teórica de fatores de risco, alocação e gestão de portfólio.
\end{enumerate}

\subsection{Recursos online}

\begin{itemize}[leftmargin=*]
    \item \textbf{QuantConnect / Lean} — plataforma open-source de backtesting e trading algorítmico.
    \item \textbf{Kaggle} — datasets financeiros e competições de ML aplicado.
    \item \textbf{Documentação oficial:} \texttt{pandas}, \texttt{scikit-learn}, \texttt{backtrader}, \texttt{statsmodels}.
    \item \textbf{YouTube:} canais de finanças quantitativas (QuantPy, Sentdex, The Trading Scientist).
\end{itemize}

% ============================================
\section{Alinhamento com Challenges}
% ============================================

O cronograma do Núcleo Quant foi desenhado para \textbf{maximizar a preparação para o Itaú Quant Challenge} e competições similares. A progressão dos blocos espelha a estrutura típica de um desafio quant:

\begin{center}
\begin{tikzpicture}[
    node distance=2cm,
    block/.style={rectangle, draw=ufuazul, fill=azulclaro, text=ufuazul, 
                  text width=3.2cm, minimum height=1.2cm, align=center, 
                  font=\small\bfseries, rounded corners=4pt}
]
    \node[block] (b1) {Bloco 1\\{\footnotesize Coleta + Análise\\de Dados}};
    \node[block, right=of b1] (b2) {Bloco 2\\{\footnotesize Risco +\\Backtesting}};
    \node[block, right=of b2] (b3) {Bloco 3\\{\footnotesize ML + Bot +\\Projeto Final}};
    
    \draw[->, thick, ufuazul] (b1) -- (b2);
    \draw[->, thick, ufuazul] (b2) -- (b3);
    
    \node[below=0.8cm of b1, font=\scriptsize\itshape, text=cinzaescuro] {Semanas 1--4};
    \node[below=0.8cm of b2, font=\scriptsize\itshape, text=cinzaescuro] {Semanas 5--8};
    \node[below=0.8cm of b3, font=\scriptsize\itshape, text=cinzaescuro] {Semanas 9--12};
\end{tikzpicture}
\end{center}

\vspace{0.3cm}

O benchmarking com o \textbf{Poli Quant} (vencedores do último Itaú Quant Challenge) já está em andamento para alinhar stack técnico, dinâmica de treinos e preparação específica. Membros do núcleo que participarem de challenges terão \textbf{prioridade de suporte} e poderão utilizar o conteúdo do núcleo diretamente nas entregas competitivas.

Além do Itaú Quant, o conteúdo do núcleo também prepara para:
\begin{itemize}[leftmargin=*]
    \item \textbf{Top Gestor} — módulo de Gestão de Risco e Otimização de Portfólio;
    \item \textbf{Hackathons de dados financeiros} — habilidades de ML e feature engineering;
    \item \textbf{Processos seletivos em buy-side e mesas quant} — domínio de Python, estatística e backtesting.
\end{itemize}

% ============================================
\section*{Considerações Finais}
% ============================================

O Núcleo Quant nasce para dar à Liga uma \textbf{vantagem competitiva real}: membros preparados para challenges quantitativos, com processo claro, entregas padronizadas e stack técnico consistente. O cronograma de 12 semanas equilibra teoria e prática, e as entregas parciais garantem que o aprendizado é contínuo — não concentrado na véspera da competição.

A meta é simples: \textbf{quem passar pelo núcleo sai sabendo montar, testar e defender uma estratégia quantitativa do zero.}

\vspace{1.5cm}

\begin{flushright}
    Atenciosamente,\\[0.5cm]
    \textbf{Arthur Pires Lopes}\\
    Diretor de Pesquisa — Núcleo Quant\\
    Liga de Mercado Financeiro — UFU
\end{flushright}

\end{document}
